本章在常背景规范下建立 Jost 解析列和散射数据, 并由传播 Lax 方程确定其演化及共振相容条件.

\section{Volterra 方程与 Jost 解析性}
固定 $z$, 记 $\Delta A=A(u)-A(Q)$. 空间 Jost 矩阵按
\[
\chi_\pm(x,w)e^{\ii\rho x\sg}\longrightarrow E(w),
\qquad x\longrightarrow\pm\infty
\]
归一化, 并令 $\mu_\pm=\chi_\pm e^{\ii\rho x\sg}$.

\begin{theorem}[Jost 解的存在性与解析列]\label{thm:jost}
设 $u-Q\in L^1(\mathbb R)\cap L^\infty(\mathbb R)$. 对连续谱上 $w\ne0,\pm\ii Q$ 的点, 存在唯一关于 $x$ 有界的归一化矩阵 $\mu_\pm$, 满足
\begin{equation}\label{eq:volterra}
\mu_\pm(x,w)=E(w)+\int_{\pm\infty}^{x}
E(w)e^{-\ii\rho(x-y)\sg}E(w)^{-1}
\Delta A(y)\mu_\pm(y,w)e^{\ii\rho(x-y)\sg}\,\dd y.
\end{equation}
列 $\mu_{-,1},\mu_{+,2}$ 在 $D_+$ 解析, 列 $\mu_{+,1},\mu_{-,2}$ 在 $D_-$ 解析; 各列在相应闭域避开 $0,\pm\ii Q$ 的紧子集上连续.
\end{theorem}
\begin{proof}
令 $W_\pm=E^{-1}\chi_\pm$, 则
\[
(W_\pm)_x=-\ii\rho\sg W_\pm+E^{-1}\Delta A\chi_\pm.
\]
在有限起点应用变分常数公式, 再利用归一化令起点趋于 $\pm\infty$, 得到
\[
\chi_\pm(x)=E e^{-\ii\rho x\sg}
+\int_{\pm\infty}^{x}E e^{-\ii\rho(x-y)\sg}E^{-1}\Delta A(y)\chi_\pm(y)\,\dd y.
\]
右乘 $e^{\ii\rho x\sg}$ 即为 \eqref{eq:volterra}.

令 $s_1=1,s_2=-1$, 取式 \eqref{eq:volterra} 第 $j$ 列, 核为
\begin{equation}\label{eq:columnkernel}
K_j(x-y,w)=
e^{\ii s_j\rho(x-y)}
E e^{-\ii\rho(x-y)\sg}E^{-1}.
\end{equation}
对 $j=1$, 中间的对角指数变为
$\diag(1,e^{2\ii\rho(x-y)})$. 当 $y\le x$ 且 $\operatorname{Im}\rho>0$ 时其模不超过 1, 故负端的第 1 列在 $D_+$ 有界. 对 $j=2$, 指数为
$\diag(e^{-2\ii\rho(x-y)},1)$; 当 $y\ge x$ 且 $\operatorname{Im}\rho>0$ 时也有界, 给出正端第 2 列. 交换符号即可得到 $D_-$ 的另两列.

取避开 $0,\pm\ii Q$ 的紧谱集 $\mathcal K$. $E,E^{-1}$ 在其上有界, 故存在 $C_{\mathcal K}$ 使允许积分方向上的 $\|K_j\|\le C_{\mathcal K}$. 例如负端第 1 列的第 $n$ 次迭代项不超过
\[
\|E_1\|\dfrac{C_{\mathcal K}^n}{n!}
\left(\int_{-\infty}^{x}\|\Delta A(y)\|\,\dd y\right)^n.
\]
阶乘来自有序积分域 $y_n<\cdots<y_1<x$. 因此 Neumann 级数绝对、局部一致收敛; 两个解之差同样迭代后由阶乘控制趋于零, 得到唯一性. 每项对 $w$ 解析, 紧集一致收敛给出解析性; 在闭域利用同一控制函数及支配收敛得到连续性. 正端积分改用反向有序域, 估计相同.
\end{proof}
在阈值 $w=\pm\ii Q$, 背景传播子为 $I+(x-y)U_\infty$, 具有线性增长. 以下解析性和散射关系均在非阈值谱上讨论.

\section{散射矩阵及谱对称}
两组空间基本解在普通连续谱上定义散射矩阵
\begin{equation}\label{eq:Sdef}
\chi_-=\chi_+S_{\rm sp}.
\end{equation}
为与传播相位一致, 在共振点之外记
\[
\Psi_\pm=\chi_\pm e^{\ii\rho Fz\sg},\qquad
\Psi_-=\Psi_+S,\qquad \mu_\pm=\Psi_\pm e^{\ii\theta\sg}.
\]
这里 $S=(s_{jk})$; $\Psi_\pm$ 的传播方程由命题 \ref{prop:evolution} 给出.

\begin{proposition}[散射系数与约化]\label{prop:scattering}
在定理 \ref{thm:jost} 的条件下,
\begin{equation}\label{eq:detjost}
\det\chi_\pm=\det\mu_\pm=\gamma,\qquad \det S=\det S_{\rm sp}=1.
\end{equation}
散射系数的 Wronskian 表达为
\begin{align}
s_{11}&=\dfrac{\det(\Psi_{-,1},\Psi_{+,2})}{\gamma},&
s_{22}&=\dfrac{\det(\Psi_{+,1},\Psi_{-,2})}{\gamma},\label{eq:wronskian}\\
s_{21}&=\dfrac{\det(\Psi_{+,1},\Psi_{-,1})}{\gamma},&
s_{12}&=\dfrac{\det(\Psi_{-,2},\Psi_{+,2})}{\gamma}.\nonumber
\end{align}
其中 $s_{11}$ 在 $D_+$ 解析, $s_{22}$ 在 $D_-$ 解析. 令
$K=\left(\begin{smallmatrix}0&1\\-1&0\end{smallmatrix}\right)$、
$C(w)=\dfrac{\ii w}{Q}\sigma_1$, 则
\begin{equation}\label{eq:conjugation}
\Psi_\pm(w)=K\conj{\Psi_\pm(\conj w)}K^{-1},\quad
\mu_\pm(w)=K\conj{\mu_\pm(\conj w)}K^{-1},\quad
S(w)=K\conj{S(\conj w)}K^{-1}.
\end{equation}
\begin{equation}\label{eq:musheet}
\Psi_\pm(h(w))=\Psi_\pm(w)C(w),\quad
\mu_\pm(h(w))=\mu_\pm(w)C(w),\quad
S(h(w))=\sigma_1S(w)\sigma_1.
\end{equation}
从而
\begin{equation}\label{eq:sconj}
s_{22}(w)=\conj{s_{11}(\conj w)},\qquad
s_{12}(w)=-\conj{s_{21}(\conj w)}.
\end{equation}
在实谱上还有 $S^\dagger S=I$.
\end{proposition}
\begin{proof}
由 $\tr\widehat U=0$, Jost 矩阵的行列式与 $x$ 无关; 在归一化端取极限得到 \eqref{eq:detjost}. Cramer 法则给出 \eqref{eq:wronskian}. 前两式中的背景指数相消, 故透射系数的解析性来自定理 \ref{thm:jost} 的两组解析列.

矩阵恒等式
\[
\widehat U(w)=K\conj{\widehat U(\conj w)}K^{-1},\qquad
E(w)=K\conj{E(\conj w)}K^{-1}
\]
以及 $\theta(w)=\conj{\theta(\conj w)}$ 表明, 共轭变换后的 Jost 列满足同一积分方程和归一化. 由唯一性得到 \eqref{eq:conjugation}, 逐元展开即为 \eqref{eq:sconj}.

另一方面,
\begin{equation}\label{eq:Esheet}
E(h(w))=E(w)C(w),\qquad C\sg=-\sg C.
\end{equation}
结合 $\lambda(h(w))=\lambda(w)$、$\theta(h(w))=-\theta(w)$, 再次由归一化唯一性得到 Jost 换叶关系, 代入散射关系即得 $S(h(w))=\sigma_1S(w)\sigma_1$. 实谱上的 $\widehat U$ 反 Hermitian, 故 $\chi_\pm^\dagger\chi_\pm$ 与 $x$ 无关, 其端值为 $E^\dagger E=\gamma I$. 因此 $S_{\rm sp}$ 及其对角酉共轭 $S$ 均为酉矩阵.
\end{proof}

在透射系数非零的轮廓部分定义反射系数
\begin{equation}\label{eq:reflection}
r=\dfrac{s_{21}}{s_{11}},\qquad
\widetilde r=\dfrac{s_{12}}{s_{22}},
\end{equation}
透射系数按相应解析域延拓, 非对角散射系数在连续谱上取值.

\section{渐近归一化与离散谱}
RHP 的两处归一化由 Jost 解析列的高能展开确定. 设
\begin{equation}\label{eq:highenergy-assumption}
u-Q\in W^{2,1}(\mathbb R).
\end{equation}
该条件保证 $u$ 有界、$u_x$ 连续且在无穷远趋于零, 并使下面分部积分中的导数可积.
\begin{proposition}[一致高能展开]\label{prop:highenergy}
在式 \eqref{eq:highenergy-assumption} 下, 令
\[
B(u)=\begin{pmatrix}0&u\\\conj u&0\end{pmatrix},\qquad
J_\pm(x)=\int_{\pm\infty}^{x}(Q^2-|u(y)|^2)\,\dd y.
\]
各 Jost 解析列在其相应闭域中满足
\begin{equation}\label{eq:mu-highenergy}
\mu_\pm(x,w)
=I+\dfrac{-\ii B(u(x))+\ii J_\pm(x)\sg}{w}
+O(w^{-2}),\qquad w\to\infty.
\end{equation}
矩阵式应按解析列理解, 余项对实 $x$ 一致. 透射系数满足
\begin{equation}\label{eq:transmission-highenergy}
s_{11}(w)=1+\dfrac{\ii J_{\rm tot}}{w}+O(w^{-2}),
\qquad
J_{\rm tot}=\int_{-\infty}^{+\infty}(Q^2-|u|^2)\,\dd x,
\end{equation}
其中 $s_{11}$ 的展开位于 $D_+$; $s_{22}$ 的展开由共轭得到.

在 $w=0$ 附近, 各列及实轴上的散射矩阵还满足
\begin{equation}\label{eq:smallmu}
\mu_\pm(w)=-\dfrac{\ii Q}{w}\sigma_1+O(1),\qquad
S(w)=I+O(w).
\end{equation}
\end{proposition}
\begin{proof}
作局部矩阵变换
\[
T(x,w)=I-\dfrac{\ii B(u(x))}{w},\qquad
\mu_\pm=Tn_\pm.
\]
利用 $A(u)=\sg B(u)$、$B(u)\sg=-\sg B(u)$、$B(u)^2=|u|^2I$. 因而
\[
\det T=1+\dfrac{|u|^2}{w^2},\qquad
T^{-1}=\dfrac{I+\dfrac{\ii B(u)}{w}}{1+\dfrac{|u|^2}{w^2}}.
\]
当 $|w|$ 足够大时, $T,T^{-1}$ 对所有实 $x$ 一致有界. 并且 $x\to\pm\infty$ 时 $T\to E$, 所以变换后 $n_\pm\to I$.

由 $\mu_x=\widehat U\mu+\ii\rho\mu\sg$ 以及 $\lambda+\rho=w$、$\rho-\lambda=\dfrac{Q^2}{w}$, 逐项计算
\begin{align*}
\widehat UT+\ii\rho T\sg-T_x
&=\ii(\rho-\lambda)\sg+A
-\dfrac{\lambda+\rho}{w}\sg B
-\dfrac{\ii}{w}AB+\dfrac{\ii}{w}B_x\\
&=\dfrac{\ii}{w}\bigl((Q^2-|u|^2)\sg+B_x\bigr).
\end{align*}
于是
\begin{equation}\label{eq:n-volterra-system}
(n_\pm)_x=-\ii\rho[\sg,n_\pm]+\mathcal C(x,w)n_\pm,
\qquad
\mathcal C=\dfrac{1}{w}T^{-1}\mathcal C_1,
\quad
\mathcal C_1=\ii\bigl((Q^2-|u|^2)\sg+B_x\bigr).
\end{equation}
由式 \eqref{eq:highenergy-assumption} 及 $u$ 有界, $Q^2-|u|^2$、$u_x$ 均可积. $T^{-1}=I+O(w^{-1})$ 一致成立, 所以
\[
\|\mathcal C(\cdot,w)\|_{L^1}\le\dfrac{C}{|w|},\qquad
\left\|\mathcal C(\cdot,w)-\dfrac{\mathcal C_1}{w}\right\|_{L^1}
\le\dfrac{C}{|w|^2}.
\]
这里和下文的 $C$ 与 $x,w$ 无关, 但可依赖势的范数. 对 $n_\pm$ 的第 $j$ 列, 积分核不再含 $E$, 而是
\[
\exp\bigl(-\ii\rho(x-y)(\sg-s_jI)\bigr).
\]
在相应解析域及积分方向中, 它的模有一致界. 由与定理 \ref{thm:jost} 相同的有序积分估计, 第 $l$ 次迭代项不超过
$\dfrac{C^l}{l!|w|^l}$. 因此 $n_\pm-I=O(w^{-1})$, 且二次及以上迭代的总和为 $O(w^{-2})$, 对 $x$ 一致.

只需计算第一次迭代. 以负端第一列为例, 其对角分量为
\[
\dfrac{1}{w}\int_{-\infty}^{x}(\mathcal C_1)_{11}(y)\,\dd y
=\dfrac{\ii J_-(x)}{w},
\]
而非对角分量为
\[
\dfrac{\ii}{w}\int_{-\infty}^{x}
e^{2\ii\rho(x-y)}\conj{u_y(y)}\,\dd y.
\]
对后一积分分部积分, 得到
\begin{align*}
\int_{-\infty}^{x}e^{2\ii\rho(x-y)}\conj{u_y(y)}\,\dd y
&=-\dfrac{\conj{u_x(x)}}{2\ii\rho}
+\dfrac{1}{2\ii\rho}\int_{-\infty}^{x}
e^{2\ii\rho(x-y)}\conj{u_{yy}(y)}\,\dd y.
\end{align*}
无穷远边界项因 $u_x\to0$ 消失. 由于指数模不超过 1, $u_x$ 有界、$u_{xx}$ 可积, 且 $|\rho|\ge c|w|$, 上式为 $O(w^{-1})$. 再乘前面的 $w^{-1}$, 非对角分量就是 $O(w^{-2})$. 正端积分和第二列只改变积分方向或指数符号, 同样给出
\[
n_\pm=I+\dfrac{\ii J_\pm\sg}{w}+O(w^{-2}).
\]
与 $T=I-\dfrac{\ii B}{w}$ 相乘即得式 \eqref{eq:mu-highenergy}.

最后, 由去振荡 Wronskian
$s_{11}=\dfrac{\det(\mu_{-,1},\mu_{+,2})}{\gamma}$, 其分子为
$1+\dfrac{\ii(J_--J_+)}{w}+O(w^{-2})$. 而
\[
J_-(x)-J_+(x)=J_{\rm tot},\qquad
\gamma=1+O(w^{-2}),
\]
故得到式 \eqref{eq:transmission-highenergy}.

在实轴的无穷远处, 将两组 Jost 展开代入散射关系得 $S=I+O(w^{-1})$; 相位因子在实谱上模为 1. 再由换叶关系
$\mu_\pm(h(w))=\mu_\pm(w)\dfrac{\ii w}{Q}\sigma_1$ 和
$S(h(w))=\sigma_1S(w)\sigma_1$, 得到 \eqref{eq:smallmu}.
\end{proof}

\begin{proposition}[简单离散谱的轨道]\label{prop:quartet}
设 $\zeta\in D_+$ 为 $s_{11}$ 的简单零点, 且远离连续谱、$0$ 和共振点. 在上述约化下,
\begin{equation}\label{eq:quartet}
\left\{\zeta,\conj\zeta,-\dfrac{Q^2}{\zeta},
-\dfrac{Q^2}{\conj\zeta}\right\}
\end{equation}
构成完整谱轨道. 其中 $\zeta$ 与 $h(\conj\zeta)$ 属于 $D_+$, 另外两点属于 $D_-$.
\end{proposition}
\begin{proof}
共轭改变 $\operatorname{Im}\rho$ 的符号, 换叶也改变其符号, 二者复合保留区域. 由 \eqref{eq:sconj}, $s_{22}(\conj\zeta)=0$; 由 \eqref{eq:musheet}, $s_{22}(h(\zeta))=0$, $s_{11}(h(\conj\zeta))=0$. 因 $h'(w)=\dfrac{Q^2}{w^2}\ne0$, 简单性保持. 离开实轴和圆周保证四点互异. Wronskian 为零说明 $\Psi_{-,1}$ 与 $\Psi_{+,2}$ 线性相关; 两列分别在左、右端衰减, 其相同向量给出平方可积本征函数.
\end{proof}

在每个简单零点定义列相关常数和范数常数
\[
\Psi_{-,1}(\zeta_j)=b_j\Psi_{+,2}(\zeta_j),\qquad
c_j=\dfrac{b_j}{s_{11}'(\zeta_j)}.
\]
去振荡后, 对应的留数因子为 $C_j=c_j e^{2\ii\theta(\zeta_j)}$. 其传播性质及配对关系分别由命题 \ref{prop:evolution} 和定理 \ref{thm:RHP} 确定.

\section{散射数据的传播演化}
设三场为包含 $z=0$ 的开区间 $I$ 上的经典 HMB 解. 假设
\[
\begin{gathered}
u-Q\in C(I;W^{3,1}(\mathbb R))\cap C^1(I;L^1(\mathbb R)),\\
v+\dfrac{\ii\eta_0Q}{a},\ \eta-\eta_0\in C(I;W^{1,1}(\mathbb R)).
\end{gathered}
\]
于是 Jost 积分方程可对 $z$ 求导, 且 $\widehat V-\widehat V_\infty$ 在两端趋于零.

\begin{proposition}[散射数据的演化]\label{prop:evolution}
在上述条件下, 对共振点之外的谱参数, $\Psi_\pm$ 同时满足两部分 Lax 方程, $S$ 与 $x,z$ 均无关, 且
\begin{equation}\label{eq:timeS}
S_{\rm sp}(w,z)=e^{\ii\rho Fz\sg}S(w)e^{-\ii\rho Fz\sg}.
\end{equation}
\begin{equation}\label{eq:timer}
r_{\rm sp}(w,z)=e^{-2\ii\rho Fz}r_{\rm sp}(w,0),\qquad
\widetilde r_{\rm sp}(w,z)=e^{2\ii\rho Fz}\widetilde r_{\rm sp}(w,0).
\end{equation}
透射系数及其离散零点不随 $z$ 改变, 同时归一化的离散列相关常数也不随 $z$ 改变.
\end{proposition}
\begin{proof}
零曲率条件给出
\[
\partial_x(\chi_{\pm,z}-\widehat V\chi_\pm)
=\widehat U(\chi_{\pm,z}-\widehat V\chi_\pm).
\]
在普通连续谱上, 由 $\widehat V_\infty E=\ii\rho F E\sg$, 左侧括号内的空间解在归一化端与
$-\chi_\pm\ii\rho F\sg$ 具有相同渐近行为. Jost 唯一性因此给出
\[
\chi_{\pm,z}=\widehat V\chi_\pm-\chi_\pm\ii\rho F\sg.
\]
右乘 $e^{\ii\rho Fz\sg}$ 即得 $\Psi_{\pm,z}=\widehat V\Psi_\pm$, 相应列恒等式延拓至各自解析域. 对散射关系求导得 $S_x=S_z=0$, 对离散点处的列相关关系求导同样得到相关常数不变. 将 $\Psi_\pm$ 的定义代回 \eqref{eq:Sdef} 得 \eqref{eq:timeS}, 其对角元不变, 非对角元分别乘 $e^{\mp2\ii\rho Fz}$, 从而得到 \eqref{eq:timer}.
\end{proof}

\section{共振点的相容性与正则性}
空间谱问题在 $\lambda=a$ 正则. 介质方程将此处的散射数据与两端背景联系起来.
\begin{proposition}[固定传播截面的共振反射条件]\label{prop:resonance}
固定一个 $z$. 设去相位后的三个场满足介质方程, 并在两端趋向 $(Q,-\dfrac{\ii\eta_0Q}{a},\eta_0)$, 其中 $a\ne0,\eta_0\ne0$. 再设空间 Jost 基底 $\chi_\pm(x,w_a^\pm)$ 正则可逆, 具有前述振荡归一化, 且 $R-R_\infty\to0$. 则在两共振原像处, 空间散射矩阵为对角矩阵; 特别地 $r_{\rm sp}=0$.
\end{proposition}
\begin{proof}
由介质方程,
\[
R_x=[\widehat U(a),R].
\]
对任一可逆空间基本解 $\chi$,
\begin{align*}
\partial_x(\chi^{-1}R\chi)
&=-\chi^{-1}\widehat U(a)R\chi
+\chi^{-1}R_x\chi+\chi^{-1}R\widehat U(a)\chi\\
&=\chi^{-1}(R_x-[\widehat U(a),R])\chi=0.
\end{align*}
故 $C_\pm=\chi_\pm^{-1}R\chi_\pm$ 与 $x$ 无关. 由 \eqref{eq:Rbg} 和 \eqref{eq:diagE},
\[
E(w_a)^{-1}R_\infty E(w_a)=
\dfrac{\eta_0\rho(w_a)}{a}\sg=:d_a\sg.
\]
在实轴 $\rho(w_a)$ 为非零实数, 振荡对角矩阵有界. 对 $C_\pm$ 分别取其归一化端点的极限, 得 $C_+=C_-=d_a\sg$. 再由 $\chi_-=\chi_+S_{\rm sp}$,
\[
d_a\sg=S_{\rm sp}^{-1}d_a\sg S_{\rm sp}.
\]
$d_a\ne0$, 故 $[\sg,S_{\rm sp}]=0$, 其两非对角元为零. 行列式为 1 保证两个对角元均非零, 反射系数因而有定义且等于零.
\end{proof}

反射系数在共振点的消失阶决定传播后的局部正则性. 对式 \eqref{eq:timer} 中的奇异相位, 有如下判据.

\begin{lemma}[有限阶零点与奇异相位]\label{prop:oscillatory-regularity}
设 $m,k$ 为非负整数, $c\in\mathbb R\setminus\{0\}$, $g$ 在零点邻域光滑, 且 $g(0)\ne0$. 对 $s\ne0$ 定义
\[
f(s)=s^m g(s)e^{\frac{\ii c}{s}}.
\]
则 $f$ 能在 $s=0$ 延拓为 $C^k$ 函数, 当且仅当 $m\ge2k+1$. 满足此条件时, 延拓的各阶导数 $f^{(j)}(0)=0$, $0\le j\le k$.
\end{lemma}
\begin{proof}
对 $j\ge0$, 在去心邻域逐次求导可写成
\begin{equation}\label{eq:oscillatory-derivatives}
f^{(j)}(s)=e^{\frac{\ii c}{s}}s^{m-2j}G_j(s),
\qquad G_j(0)=(-\ii c)^j g(0),
\end{equation}
其中 $G_j$ 光滑. 对 $j=0$ 取 $G_0=g$. 若式 \eqref{eq:oscillatory-derivatives} 对 $j$ 成立, 由
\[
\dfrac{\dd}{\dd s}e^{\frac{\ii c}{s}}
=-\dfrac{\ii c}{s^2}e^{\frac{\ii c}{s}}
\]
可得
\[
G_{j+1}(s)
=-\ii cG_j(s)+(m-2j)sG_j(s)+s^2G_j'(s).
\]
这一递推同时证明光滑性及零点处的取值公式.

若 $m\ge2k+1$, 对 $0\le j\le k$ 有 $m-2j\ge1$, 故各个去心导数都趋于零. 还需确认这些极限确实是延拓函数的导数. 对 $j<k$, 假定第 $j$ 阶延拓已定义, 则
\[
\dfrac{f^{(j)}(s)-f^{(j)}(0)}{s}
=e^{\frac{\ii c}{s}}s^{m-2j-1}G_j(s)\longrightarrow0,
\]
因为 $m-2j-1\ge2$. 从 $j=0$ 起逐阶应用此式, 即得 $C^k$ 延拓且各阶导数在零点为零.

反之, 若 $m<2k$, 式 \eqref{eq:oscillatory-derivatives} 的 $j=k$ 项因 $G_k(0)\ne0$ 而在模上无界; 若 $m=2k$, 该项包含振幅趋于非零常数的 $e^{\frac{\ii c}{s}}$, 没有极限. 两种情形均不可能为 $C^k$.
\end{proof}

\begin{proposition}[共振传播的局部正则性]\label{prop:resonant-regularity}
在 $\lambda=a$ 附近固定一张实谱叶, 以 $\lambda$ 为局部坐标表示 $r_{\rm sp}$, 并设其初值具有有限阶零点
\[
r_{\rm sp}(a+s,0)=s^m g_0(s),\qquad g_0\in C^\infty,\quad g_0(0)\ne0,
\]
其中 $m$ 为非负整数. 对任意固定 $z\ne0$ 和非负整数 $k$, 演化式 \eqref{eq:timer} 给出的反射系数可在 $s=0$ 延拓为 $C^k$ 函数, 当且仅当 $m\ge2k+1$. 该判据在两个单值化共振原像均成立.
\end{proposition}
\begin{proof}
记 $\rho_a=\rho(a)=\pm\sqrt{a^2+Q^2}$, 并写
$F_{\rm pol}(\lambda)=4\beta\lambda^2-\alpha\lambda-2\beta Q^2$.
由于 $\rho$ 在所选实谱叶上解析, 函数
\[
L(s)=\rho(a+s)F_{\rm pol}(a+s)
+\dfrac{\eta_0}{a}\dfrac{\rho(a+s)-\rho_a}{s}
\]
在 $s=0$ 可去. 因此
\begin{equation}\label{eq:local-resonant-phase}
-2z\rho(\lambda)F(\lambda)=\dfrac{c_z}{s}-2zL(s),\qquad
c_z=-\dfrac{2z\eta_0\rho_a}{a}\ne0.
\end{equation}
演化后的反射可写为
\[
r_{\rm sp}(a+s,z)=s^m g_z(s)e^{\frac{\ii c_z}{s}},\qquad
g_z(s)=g_0(s)e^{-2\ii zL(s)}.
\]
$g_z$ 光滑且 $g_z(0)\ne0$, 应用引理 \ref{prop:oscillatory-regularity} 即得充要条件. 又
\[
\lambda'(w_a^\pm)=\dfrac{1}{2}\left(1+\dfrac{Q^2}{(w_a^\pm)^2}\right)>0,
\]
故 $\lambda-a$ 与 $w-w_a^\pm$ 为等价的局部光滑坐标, 消失阶和 $C^k$ 性质不变.
\end{proof}
